% TOP_PROBLEM: {"id": 3039, "title": "Rainbow AP(4) in an almost equinumerous coloring", "category_id": 2, "category": {"id": 2, "name": "combinatorics", "display_name": "Combinatorics", "description": "Counting problems, graph theory, discrete structures.", "slug": "combinatorics", "order_index": 2, "created_at": "2026-07-31T15:26:25.671Z"}, "status": "open", "problem_number": "OPG-478", "set_id": 12}
% FIRST_POSTED: 2026-10-01
% ATTEMPT_STATUS: unresolved
% UnsolvedMath ID 3039; source catalogue code OPG-478
% !TeX program = lualatex
% Research attempt by GPT-6 Astra Ultra; no claim of a new complete solution.
\documentclass[11pt,a4paper]{article}
\usepackage[margin=25mm]{geometry}
\usepackage{fontspec}
\setmainfont{lmroman10-regular.otf}[BoldFont=lmroman10-bold.otf,
 ItalicFont=lmroman10-italic.otf,BoldItalicFont=lmroman10-bolditalic.otf]
\usepackage{amsmath,amssymb,mathtools}
\usepackage{unicode-math}
\setmathfont{latinmodern-math.otf}
\AtBeginDocument{\let\setminus\smallsetminus}
\usepackage{xurl,hyperref}
\hypersetup{colorlinks=true,urlcolor=blue}
\setcounter{secnumdepth}{0}
\setlength{\parindent}{0pt}
\setlength{\parskip}{5pt}
\setlength{\emergencystretch}{3em}
\begin{document}
\section*{Rainbow AP(4) in an almost equinumerous coloring}\label{3039}
{\small UnsolvedMath ID 3039 · OPG-478 · Combinatorics · OpenGarden}
\subsection{Definitions and mathematical statement}
For a prime $p\ge5$, let $\mathbb Z_p=\mathbb Z/p\mathbb Z$ with addition
modulo $p$. A four-colouring is a map $\chi:\mathbb Z_p\to\{1,2,3,4\}$,
with colour classes $C_i=\chi^{-1}(i)$. Call it almost equinumerous if
\[
 |C_i|\in\{\lfloor p/4\rfloor,\lceil p/4\rceil\}
 \qquad(1\le i\le4).
\]
Thus all classes differ in size by at most one; the sizes necessarily sum
to $p$. This is a condition on all residues, not on intervals of representatives.

A four-term arithmetic progression in $\mathbb Z_p$ is an ordered list
$a,a+d,a+2d,a+3d$ with $a\in\mathbb Z_p$ and
$d\in\mathbb Z_p\setminus\{0\}$. Its four entries are distinct because
$p\ge5$. It is rainbow when the four values of $\chi$ on these entries are
pairwise different.

Does there exist an integer $p_0$ such that for every prime $p\ge p_0$ and
every almost equinumerous four-colouring, a rainbow four-term arithmetic
progression exists? The threshold is independent of the colouring. A
negative answer requires counterexamples for arbitrarily large primes,
not merely an exceptional small prime. Progressions may wrap around the
chosen residue representatives, since the equations are in $\mathbb Z_p$.
The source asks for the prime cyclic group: substituting a composite-order
cyclic group changes the problem. No particular common difference or
particular order of the four colours is prescribed.
\subsection{Short English statement}
Must every sufficiently large prime cyclic group, split into four almost equal colour classes, contain a four-term progression using all four colours?
\subsection{Research attempt}
\paragraph{Scope and process.}
This is a GPT-6 Astra Ultra research attempt dated 1 October 2026, using
primary-source browsing and elementary mathematical checks. The canonical
definition above is retained. Results below are restricted results or
reductions, not a claim of a new solution to the entire catalogue question.
No formal proof assistant or external mathematical referee checked this text.


\paragraph{Literature and chosen attack.}
The original page [S2, R1] explicitly separates prime cyclic groups
from composite cyclic groups, where balanced rainbow-free examples
are possible. We keep the prime hypothesis and examine two approaches:
a constructive interval subcase, and counting monochromatic pairs
inside progressions. The second approach exposes a quantitative deficit
rather than providing a universal proof.

\paragraph{A proved subcase: four consecutive arcs.}
Suppose each colour class is one interval of consecutive residues around
the circle. Rotate the circle and rename colours so these intervals
occur in order $C_0,C_1,C_2,C_3$, starting at zero. Write
$p=4q+r$ with $0\le r\le3$. Each interval has length $q$ or $q+1$,
and its left endpoint is $b_j=jq+s_j$, where $0\le s_j\le3$ counts
the preceding long intervals. If $q\ge4$, choose any integer
$a\in[3,q-1]$. Then for $j=0,1,2,3$,
\[
 b_j\le a+jq\le b_j+q-1.
\]
The first inequality uses $s_j\le3\le a$, and the second uses
$a\le q-1\le q-1+s_j$. Thus $a+jq\in C_j$, and the progression
with difference $q$ is rainbow. Its last term is at most $4q-1<p$,
so no wraparound convention is being hidden. There are at least $q-3$
such ordered progressions with this difference. Translating back proves
the assertion for every balanced four-arc colouring when $p\ge17$.
No primality was actually needed for this restricted construction.

\paragraph{Exact collision count for an arbitrary colouring.}
Let $n_i=|C_i|$ and, for $(a,d)\in\mathbb Z_p\times\mathbb Z_p^*$,
let $Q(a,d)$ count the unordered pairs of positions in
$(a,a+d,a+2d,a+3d)$ with equal colours. For each $0\le j<k\le3$,
the map
\[
 (a,d)\longmapsto(a+jd,a+kd)
\]
is a bijection onto ordered pairs of distinct residues: recover
$d=(y-x)/(k-j)$ and $a=x-jd$. This uses $p\ge5$, so $k-j$ is
invertible. Hence
\[
 \sum_{a,d\ne0}Q(a,d)=6\sum_{i=1}^4n_i(n_i-1).
\]
Every non-rainbow progression has $Q\ge1$. If $R$ is the number of
rainbow ordered progressions, the union bound therefore yields
\[
 R\ge p(p-1)-6\sum_i n_i(n_i-1).
\]
For balanced sizes the right side is
$-p^2/2+5p+O(1)$, which is negative for large $p$. Thus this exact
two-point information is too weak to force even one rainbow progression.
Balance does not supply the higher-order correlation information missing
from this inequality.

\paragraph{Why a random calculation also does not prove the target.}
Uniformly assign the four colours with fixed positive class sizes $n_i$.
For any fixed progression its four distinct positions have four
different colours with probability
\[
 \frac{4!\,n_1n_2n_3n_4}{p(p-1)(p-2)(p-3)}.
\]
Multiplying by $p(p-1)$ gives the exact expected number of rainbow
progressions. For balanced sizes it is $(3/32+o(1))p^2$, so some
balanced colourings have many. The quantifier in the problem is every
balanced colouring, however; a positive expectation neither excludes
zero-count colourings nor establishes a uniform lower bound.

\paragraph{Verification and first remaining gap.}
The interval proof permits the long classes in any order and uses no
special prime residue class modulo four. The counting proof counts
ordered $(a,d)$, including reversals, consistently on both sides.
The missing result is a bound on structured colour correlations for
arbitrary balanced classes strong enough to force $R>0$. Neither
average colour-class sizes nor the two-point identity gives it.

\paragraph{Final status.}
\textbf{UNRESOLVED.} The consecutive-arc subcase and exact counting
identities are proved; no threshold valid for all balanced colourings
is established.

\subsection{Sources}
\begin{itemize}
\item[S1] ULAM.AI and the UnsolvedMath Contributors, supplied problems.json, record 3039 (OPG-478), OpenGarden collection. Catalogue identity and source transcription; CC BY 4.0.
\url{https://huggingface.co/datasets/ulamai/UnsolvedMath}.
\item[S2] Open Problem Garden, Rainbow AP(4) in an almost equinumerous coloring. Original problem and discussion; the mathematical formulation below is expanded from the supplied source record.
\url{http://www.openproblemgarden.org/op/rainbow_ap_4_in_an_almost_equinumerous_coloring}.

\item[R1] David Conlon, problem as posted by Veselin Jungi\'c,
\emph{Rainbow AP(4) in an almost equinumerous coloring}, original
problem page and bibliography, checked 1 October 2026.
\url{https://www.openproblemgarden.org/op/rainbow_ap_4_in_an_almost_equinumerous_coloring}.

\end{itemize}
\end{document}
